% !Mode:: "TeX:UTF-8"
% !TEX root = ../main.tex

\iffalse
\bibliography{static/references/reference.bib} % 欺骗latextools获取bib文件
\fi
\xfbody

\chapter{引言}
    \section{研究背景和意义}
           \xfbody{根据商务部发布的中国新电商发展报告(2025)\cite{spb2024_bulletin_20250522}可以看出，2024年中国网络零售额达15.5万亿元，同比增长7.2\%，拉动社会消费品零售总额增长1.7个百分点，连续12年稳居全球最大网络零售市场如图\ref{fig:kuaidi1}所示。随着电商行业的快速发展，仓储物流成为制约电商发展的重要瓶颈之一。传统的人工仓储模式效率低下，难以满足日益增长的订单需求和复杂的仓储环境。因此，智能仓储自主移动机器人作为一种新兴的仓储自动化解决方案，受到了广泛关注。它们能够在仓库中自主导航、搬运货物，提高仓储效率和准确性，降低人力成本。\par
            \begin{figure}[htbp]
                \centering
                \includegraphics[width=0.64\textwidth]{kuaidi1.png}
                \caption{2020-2024年人均用邮支出、快递支出和快递使用量情况}
                \label{fig:kuaidi1}
             \end{figure}
        
           智能仓储自主移动机器人以自动引导小车（Automated Guided Vehicle，AGV）为代表，具有自动导航、路径规划、货物识别和搬运等功能。它们能够在复杂的仓储环境中自主移动，避免障碍物，提高仓储效率和准确性。智能仓储机器人不仅能够提升仓储效率，还能够降低人力成本，减少劳动强度，提高工作安全性。如图\ref{fig:AGV}为AGV小车工作示意图。因此，研究智能仓储自主移动机器人拥有重要的理论意义和实际应用价值。通过深入研究智能仓储机器人，可以推动仓储自动化技术的发展，提高仓储效率和服务质量，促进电商行业的持续发展。同时，智能仓储机器人也可以应用于其他领域，如制造业、医疗等，具有广泛的应用前景。统计数据显示，2019年中国AGV机器人行业市场规模48.27亿元，2024年中国AGV机器人行业市场规模95.49亿元。
        \begin{figure}[htbp]
                \centering
                \includegraphics[width=0.64\textwidth]{AGV.png}
                \caption{AGV小车工作示意图}
                \label{fig:AGV}
             \end{figure}
        }
       

    \section{国内外AGV机器人研究综述}
       \xfbody{
        \subsection{AGV小车国内研究现状}
            \xfbody{我国的AGV研究起步于20世纪70年代末，并在“中国制造 2025”等政策的强力驱动下，已从最初的仿制阶段迈向自主创新的快速增长期。国内的研究主要集中在机器人导航、路径规划、货物识别和搬运等方面。自从亚马逊公司推出Kiva系统以来，我国许多企业开始模仿研发仓储机器人，形成了以AGV为主的仓储机器人产业链。国内的研究机构和企业在仓储机器人领域取得了一系列重要成果，例如京东、阿里巴巴等电商巨头都在积极研发和应用智能仓储机器人，以提升仓储效率和服务质量。\par
            AGV产业的发展不断壮大，市场需求变多，应用的行业也越来越多。实现AGV的智能化、信息化、网络化，可以让企业管理越来越方便，生产越来越安全。AGV小车具有自动化程度高、效率高、适应性强等优点，能够有效提升仓储物流的效率和准确性。随着技术的不断进步，AGV小车在仓储物流中的应用将会越来越广泛，为企业带来更多的经济效益和竞争优势。\par
            目前，国内AGV在导航技术上取得了显著突破，已实现从传统的磁条感应、二维码识别向激光SLAM和视觉SLAM等自然导航方式的全面进化。视觉SLAM因其无需铺设辅助标志、成本低和信息获取方便且丰富的优势，正逐渐成为国内达规模应用的重要研究方向。在具体的行业应用中，除汽车制造和烟草等传统领域外，电商物流、智慧医院、深海探测及特种危险环境也成为了国内AGV技术落地的重要场景。\par
            在算法研究层面，国内学者正钻研于复杂动态环境下的路径规划与多机协同调度技术，力求解决大规模AGV系统的冲突与效率问题。目前的研究热点涵盖了从改进的A*算法、Dijkstra算法到遗传算法、蚁群算法等生物启发式智能算法，并涌现出如黏菌算法（SMA）、黑猩猩优化算法等新兴方案。这些算法在提升AGV系统的路径规划效率和调度性能方面展现出显著优势，推动了国内AGV技术的持续创新与应用拓展。\par
            
            总之，国内的AGV市场正在快速发展，技术水平不断提升，应用领域也在不断扩大。随着5G、物联网和人工智能等技术的不断发展，我国的AGV小车在未来将会有更广阔的应用前景，为企业带来更多的经济效益和竞争优势。\par
            }
        \subsection{AGV小车国外研究现状}
            \xfbody{国外在机器人研究领域起步较早，凭借在控制技术、精密制造和传感器领域的深厚积累，长期处于全球技术领先地位。以德国、日本和美国为代表的传统制造强国，根据不同行业需求研发出了种类繁多的机器人系统\cite{Chand2025SmartCityMobileRobots}，不仅奠定了智能化机器人的核心技术基础，还引领了从工业应用向服务业及特种领域的全面转型。国外的研究主要集中在机器人导航、路径规划、货物识别和搬运等方面。国外的研究机构和企业在仓储机器人领域取得了一系列重要成果，例如亚马逊、谷歌等科技巨头都在积极研发和应用智能仓储机器人，以提升仓储效率和服务质量。\par
            在仓储物流领域，国外的研究与应用的商业化水平已经高度成熟。亚马逊公司于2012年收购的Kiva系统是货架到人模式的开创者，通过移动货架显著提升了拣选效率，成为全球智能仓储的引领者。此外，德国Magazino公司推出的TORU机器人找到了更先进的智能化方向，它不仅能实现自主导航，还能从货架上精确识别并拾取单个物品。针对复杂地形，国外研究涵盖了从单足到多足的多种仿生结构，其中波士顿动力（Boston Dynamics）开发的Atlas液压驱动机器人具备28个自由度，可以做出跑步、跳跃等动作，在处理突发情况具有极强的自主性，体现了目前国外在智能仿生机器人领域的先进水平。\par
            在算法研究层面，国外学者广泛应用启发式算法（如改进的A*、Dijkstra）、群智能优化算法以及强化学习来解决多机器人路径规划和任务分配。目前的研究趋势正向分布式协同、云协同和人工智能集成方向迈进，通过多Agent系统的自我演化和深度强化学习（如DQN），赋予机器人更强的群体协作能力和对动态环境的自适应能力。总之，国外机器人研究现状体现出结构拟人化、功能集成化和决策智能化的显著特点。\par
            国外对移动机器人（包括AGV）的研究起步较早，凭借在控制技术、制造工艺和传感器领域的深厚积累，长期处于全球技术领先地位。早在1959年，美国研制出了世界上第一台工业机器人Unimate，开启了工业自动化应用。随后在1966年至1972年间，斯坦福大学研发出世界上首个具有人工智能的移动机器人Shakey，为视觉感知、路径规划和动态避障等现代AGV核心技术奠定了基础发展到现在，国外 AGV 技术已相对成熟。从最初的磁导航和激光导航技术，到现在的基于机器视觉和深度学习技术的自主导航AGV，AGV技术在不断提高自主性和智能性的同时，开发的应用场景也越来越多。\par
            总之，国外在AGV小车的研究方面已经取得了显著的进展，尤其是在导航技术、路径规划算法和多机器人协同方面。随着人工智能和机器学习技术的不断发展，未来国外在AGV小车领域的研究将继续向更高水平的智能化和自动化方向发展。\par
            }
          }
     \section{国内外AGV机器人路径规划算法研究现状}
               \xfbody{在仓储机器人路径规划算法的研究方面，国内外学者都取得了显著的进展。黄晓宇，孙勇智，李津蓉，王翼挺，杨颀伟提出了一种模型预测控制（MPC）和微分先行比例-积分-微分（PID）协同的双闭环控制策略，用于麦克纳姆轮全向移动平台的轨迹跟踪控制\cite{黄晓宇2023基于}。曹梦龙，赵文彬，陈志强指出了机器人路径规划的步骤和传统算法的局限性，概述了群体智能优化算法及其在路径规划中的作用\cite{cao2023robot}。白宇鑫，陈振亚等人描述了一种改进的哈里斯鹰优化算法，用于解决机器人路径规划问题，讨论了HHO算法在收敛性和局部最优问题上的局限性\cite{白宇鑫2025基于改进哈里斯鹰算法的机器人路径规划研究}。赖荣燊，窦磊，巫志勇，孙帅指出了传统A*算法和动态窗口法各自的优缺点以及现有研究的改进方向\cite{赖荣燊2024融合改进}。张瑞，周丽，刘正洋提出了RRT*和DWA算法融合的路径规划方法，旨在解决复杂动态环境下移动机器人路径规划的安全性和最优性问题\cite{张瑞2024融合}。\par
               总之，国内外在AGV机器人路径规划算法的研究方面已经取得了显著的进展，尤其是在启发式算法、群智能优化算法和深度强化学习等方面。未来的研究将继续向更高水平的智能化和自动化方向发展，以满足日益复杂的仓储环境和多样化的应用需求。\par
               国外在AGV机器人路径规划算法方面的研究也取得了重要进展。M Rybczak, N Popowniak, A Lazarowska等人提出了移动机器人控制中复杂的机器学习算法和有限的计算资源、实时决策与运动控制、算法对环境变化的适应性、获取大量有价值数据以及确保机器人操作的安全性和可靠性等挑战\cite{rybczak2024survey}。MFR Lee, SH Yusuf为移动机器人引入了深度Q网络（DQN）和双深度Q网络（DDQN）两种深度Q学习智能体，使其能够自主学习避撞和导航能力\cite{lee2022mobile}。Y Dai, J Yu,C Zhang,B Zhan, X Zheng描述了传统的鲸鱼优化算法（WOA）及其改进方法在移动机器人路径规划中存在收敛速度慢、易陷入局部最优和缺乏动态避障能力等问题\cite{dai2023novel}。S Das, SK Mishra提出了一种基于机器学习概念的新颖方法，即自适应随机梯度下降线性回归（ASGDLR）算法，用于区分自主移动机器人（AMR）的右转和左转方向运动\cite{das2022machine}。B Tao,JH Kim提出了一种具有随机扰动策略的双种群粒子群优化（BPPSO）算法，通过将粒子分成两个子种群来应对传统粒子群优化算法在解决约束问题、容易陷入局部最优和早熟收敛方面的局限性\cite{tao2024mobile}。\par
               总之，国外在AGV小车路径规划算法的研究方面依然保持着高度的创新和活跃，在机器人的控制算法，避障算法和路径规划算法等方面不断取得新的突破，推动了AGV技术的持续发展和应用拓展。\par
               \par
               }
     \section{研究方法与手段}
             \subsection{研究内容}
                \xfbody{ 机器人本体结构也日益多样化，形态多样的AGV小车不断涌现，以适应不同的仓储环境和搬运需求。例如本文将要研究的潜伏式举升型机器人，其是一种能够实现自主载货、升降、运输的自动化物流设备。其特点是配备了举升装置，可以将货物从地面抬升到较高的位置，从而实现搬运和存储等任务。通常用于工业和仓储物流应用中，可以有效地提高生产效率和降低物流成本。举升型机器人具有高拓展性、高柔性、场景适应性强等优点。\par
                对仓储搬运机器人的底座的设计，包括底座的工作方式、电动机的选择、移动平台的设计。本机构大体分为三个部分进行设计，升降机构是完成载物台的升降工作，是产品的执行机构。行走装置是实现AGV的行走移动，是产品的移动机构。避障装置是通过视觉或者红外，超声波等检查系统，完成对AGV前方障碍物的检测，避免撞车等。\par
                选用叉剪式做本文的升降结构设计，其结构有传动稳定，动力输出小的特点广泛的应用于AGV这类小结构的设计，而桅柱式、臂架式更多的应用于大型升降结构，而套筒油缸式主要是应用于重型起重结构，而桁架式更用于工厂。\par
                将建立仓储搬运机器人的差速驱动底盘运动学模型，差速驱动底盘是移动机器人（如AGV、仓储机器人）最常见的运动控制形式之一，其核心结构通常由两个独立驱动的轮子和若干万向轮组成。该模型通过主控芯片分别向左右两个驱动轮发出不同的速度指令，从而驱动机器人完成直行、转弯或原地调头等动作。\par
                本文选用了黏菌算法作为循迹算法。黏菌算法是模拟黏菌捕食行为的一种仿生算法。黏菌根据气味识别到食物后，通过细胞质的流动形成多个叶脉，同时向多个食物进行移动，当叶脉接近食物源时，黏菌的生物振荡器产生一种传播波，增加细胞质流动速度，细胞质流动越快，叶脉就越厚。这种正反馈机制，使得黏菌能以一种较优的方式建立连接食物的路径。SMA具有收敛速度快，寻优能力强的优点。\par
                本文将选用pid控制方法，PID（比例-积分-微分）控制是机器人和自动化系统中最核心的反馈控制算法之一。它通过计算目标值与实际反馈值之间的误差，并对该误差进行三种方式的处理来实时调整控制指令。PID控制示意图如图\ref{fig:pid}所示。\par
               \begin{figure}[htbp]
                \centering
                \includegraphics[width=0.64\textwidth]{pid.png}
                \caption{PID控制示意图}
                \label{fig:pid}
             \end{figure}
     }
       
          通过上述研究，本文将系统地设计和分析智能仓储机器人，旨在提升其性能和应用效果，为智能仓储领域的发展做出贡献。\par
            \subsection{研究方法}
                \xfbody{本文将采用理论分析的研究方法。首先，通过理论分析建立智能仓储机器人的运动学模型和控制算法，深入理解其工作原理和性能特点。其次，通过仿真软件对机器人运动性能进行全面评估，分析运动轨迹、速度、加速度等关键指标，验证设计的动态性能。\par
                利用三维建模软件对机器人各部件进行精确建模，完成整机虚拟装配，通过装配配合检验结构设计的合理性与可行性，及时发现并修正设计缺陷。\par
                系统查阅国内外关于仓储搬运机器人、移动机器人底盘结构、机械臂设计、运动学建模、模型pid控制算法以及黏菌算法（SMA）等相关领域的文献资料，梳理现有研究成果与技术瓶颈，为本研究提供理论支撑与技术参考。\par
                总之，本文将通过理论分析、仿真评估、三维建模和文献综述等多种研究方法，系统地设计和分析智能仓储机器人。\par
                

                }

     \section{论文主要研究内容}   
           \xfbody{本文的主要研究内容包括以下几个方面：\par
           1. AGV总体设计方案：设计满足现代智能仓储需求的AGV小车，具备自主导航能力，能够在密集货架环境中灵活穿梭，完成物料的自动搬运与存取任务。系统采用基于差速驱动的底盘结构，结合智能避障与路径规划算法，实现高效、安全的自动行走与作业。\par
           2. AGV机械结构设计：设计AGV小车的核心部件结构，包括载物装置、行走装置和避障装置等核心模块，形成完整的智能搬运系统。进行结构强度校核和有限元分析，确保AGV在复杂环境下的稳定运行。\par
           3. AGV电气系统设计：设计AGV小车的电气系统总体架构，包括电机、传感器等部件选型与设计，以及控制电路设计，确保AGV能够实现自主导航和智能避障功能。\par
           4. AGV运动学建模与仿真：建立AGV小车的差速驱动底盘运动学模型，并搭建MATLAB/Simulink仿真模型，对典型运动轨迹进行仿真与分析，验证设计的动态性能。\par
           5. AGV运动控制算法分析：分析PID控制算法原理，设计双轮差速PID控制器，实现基本运动控制，并通过仿真验证控制效果。\par
           6. 结论：总结全文工作，分析不足和改进方向，并展望AGV小车在智能仓储领域的应用前景。\par
           
     
             }
\chapter{AGV总体设计方案}
    \section{设计需求和技术指标}
         \xfbody{为满足现代智能仓储对空间利用率和作业灵活性的要求，本文设计了一款具有自主导航能力的移动机器人。该机器人能够在密集货架环境中灵活穿梭，完成物料的自动搬运与存取任务。系统采用基于差速驱动\cite{刘梓博2025基于双变量限幅}的底盘结构，结合智能避障与路径规划算法，实现高效、安全的自动行走与作业。\par
         该AGV小车的主要技术指标如下：\par
         自动载物功能：配备可升降的载物台，升降行程为0.5m，可满足不同货架高度的物料搬运需求。配合不同的载物平台，可以适应多种类型的货物。\par
         自动循迹能力：最大外形尺寸控制在$1 \times 1 \times 1\,\mathrm{m}$，保证在狭窄通道与密集货架间的灵活通行的同时满足大量物品运输的要求\par
         结构组成：整车包括载物装置、行走装置、避障装置等核心模块，形成完整的智能搬运系统。具备多种传感器，提升了环境感知能力，提高其在复杂环境中移动的能力。\par
         载重能力：最大承载重量不超过25kg，适用于轻型物料的仓储搬运场景。 \par
         智能避障与导航：通过传感器融合技术（如激光雷达、超声波、视觉模块），实现对环境的实时感知与动态避障，保证在复杂仓储环境下的安全运行。\par
         系统优势：具备高空间利用率、灵活性强、自动化程度高，能够显著提升仓储作业效率，降低人工成本的同时提高作业质量。\par



         }
    \section{总体结构和驱动选型}
           \subsection{总体结构设计}
                \xfbody{本文设计的举升式AGV车体结构在整体方案上采用焊接钢架作为主体框架，以保证车体的整体强度和稳定性。为了在有限空间内兼顾承载能力与结构紧凑性，底盘部分选用厚度为15mm的钢板，通过焊接成型后再进行精加工处理，以确保底盘的定位精度和结构强度，从而满足AGV在复杂仓储环境下的长期稳定运行需求。\par
                在底盘上安装有万向轮与驱动件，驱动系统采用差速驱动方式，并在差速轮处配置减振机构，有效降低运行过程中因地面不平或高速行驶所产生的振动，提高车辆的平稳性和控制精度。为了实现人机交互的便利性，控制板被合理布置在车体内部，既保证了电气系统的安全性，又便于维护；同时，控制面板安装在车头位置，使操作人员能够直观地进行操作与监控。\par
                该结构设计不仅兼顾了强度、稳定性与紧凑性，还为后续的自动导航、避障与举升功能提供了坚实的硬件基础。整体方案在满足载重与尺寸要求的同时，兼顾了空间利用率与操作便利性，体现了现代智能仓储对AGV设备的高效性与灵活性需求。\par
                }
           \subsection{驱动选型}
                \xfbody{在驱动系统的选型上，本文选择了差速驱动方式,如图\ref{fig:chasu}所示。
                  \begin{figure}[htbp]
                \centering
                \includegraphics[width=0.64\textwidth]{chasu.png}
                \caption{差速驱动系统示意图}
                \label{fig:chasu}
             \end{figure}
               作为AGV小车的主要驱动形式，差速驱动系统由两个独立控制的轮子组成，通过调整左右轮的速度差来实现转向和行驶控制。这种驱动方式具有结构简单、控制灵活、适应性强等优点，特别适合在狭窄空间和复杂环境中运行的AGV小车。\par
                具体来说，差速驱动系统能够实现原地转弯、灵活转向和高效行驶，适应不同的仓储环境和作业需求。通过合理设计轮子的尺寸和驱动电机的功率，可以满足AGV小车的载重能力和速度要求。同时，差速驱动系统的控制算法相对简单，便于实现高精度的运动控制和路径规划，进一步提升AGV小车的自主导航能力和作业效率。\par
                差速驱动系统的选型不仅满足了AGV小车在仓储环境中的灵活性和适应性需求，还为后续的运动控制算法设计提供了良好的基础，确保AGV小车能够高效、安全地完成物料搬运任务。\par
                在该装置的设计中，驱动部分的壳体采用 45号钢板 制成，厚度为 2mm，通过弯曲和焊接工艺加工而成，既保证了整体结构的强度，又兼顾了重量的合理控制。驱动级的两个电机分别安装在固定壳体的左右两侧，形成对称布局，以确保传动过程中的受力均衡。电机的另一侧与驱动装置的底座相连并固定，使得整个驱动系统在运行时更加稳定。\par
                驱动装置的底座位于整体结构的下方，并与壳体紧密结合。底座顶部安装有轴承，孔内嵌入直线轴承模块，支撑杆则固定在直线轴承模块上，从而保证支撑杆在安装时能够保持上下同心，避免因偏心而导致的振动或磨损。\par
                在驱动结构中，减振器由多个关键部件组成，包括 预紧螺母、支撑立轴、车身骨架连接板、固定板、阻尼弹簧以及直线轴承板。其中，支撑立轴安装在直线轴承的中间位置，以确保其同心度。预紧螺母位于支撑立轴的顶部，用于减振器的安装与调节；支撑立轴的下端通过圆柱销与驱动装置的底座连接并固定，进一步增强了结构的稳定性。\par
                减震弹簧则压缩在车架连接板与固定板之间，形成弹性缓冲区。当驱动装置运行时，弹簧能够有效吸收和缓解传动过程中产生的冲击与振动，从而使整个传动部分达到良好的减震效果。这种结构设计不仅提升了装置的耐用性和运行平稳性，还延长了各零部件的使用寿命。\par
                以上设计方案充分考虑了AGV小车在复杂仓储环境中的运行需求，通过合理的结构布局和驱动系统选型，确保了AGV小车能够高效、安全地完成物料搬运任务，同时为后续的运动控制算法设计提供了坚实的硬件基础。\par
                }    
    \section{系统总体架构}
           \xfbody{AGV小车的系统总体架构主要包括机械结构、电气系统、运动控制算法和软件系统四个核心部分。机械结构方面，AGV小车采用差速驱动底盘设计，配备升降装置和避障装置，以满足仓储环境中的搬运需求。电气系统方面，AGV小车配备了电机、传感器和控制电路，实现自主导航和智能避障功能。运动控制算法方面，AGV小车采用PID控制算法，实现基本运动控制，并通过仿真验证控制效果。软件系统方面，AGV小车搭载了基于ROS（Robot Operating System）的操作系统，实现对硬件的统一管理和控制，同时支持路径规划、任务调度等高级功能。\par
           采用举升式AGV的方案设计，让小车潜入作业车的底部，可以实现自动挂接和分离后再进行牵引作业。与传统的AGV相比，举升式AGV具有更高的空间利用率和更强的适应性，能够在狭窄的仓储环境中灵活穿梭，提高搬运效率。同时，举升式设计还可以减少对地面环境的依赖，降低对地面平整度的要求，使得AGV能够在更多样化的仓储环境中应用。\par
           在传感器选择上，根据现代仓储物流的需求和面对仓库复杂的空间要求，对比多种传感器，本文选择了超声波传感器，红外线传感器和摄像头模块作为本AGV机器人的主要传感器。超声波传感器具有较长的测距范围和较强的抗干扰能力，适用于检测较远距离的障碍物；红外线传感器则适合近距离的障碍物检测，能够提供更高的精度；摄像头模块则可以实现视觉识别和环境感知，增强AGV的自主导航能力。这些传感器的组合能够有效提升AGV在复杂仓储环境中的避障能力和导航性能。\par
           在AGV小车机械元件上，本文选择了Q235A钢作为主要材料。Q235A钢具有良好的机械性能和较高的强度，能够满足AGV小车在仓储环境中承受的载荷和冲击要求。同时，Q235A钢具有较好的焊接性能和加工性能，便于制造和维护。此外，Q235A钢的成本相对较低，有助于控制AGV小车的制造成本，提高其市场竞争力。\par
           AGV小车的动力源选择了电池作为主要的能源供应方式。电池具有高能量密度、可充电和环保等优点，适合AGV小车在仓储环境中的使用需求。通过合理设计电池容量和充电系统，可以确保AGV小车在长时间运行过程中保持稳定的动力供应。同时，电池驱动的AGV小车还能够实现更高的效率和更低的维护成本，进一步提升AGV在智能仓储中的应用价值。电机、电池和电控系统的选型将根据AGV小车的载重能力、运行速度和续航时间等技术指标进行综合考虑，以确保AGV小车能够高效、安全地完成物料搬运任务。电机将选择伺服电机，以提供精确的速度和位置控制；电池将选择锂离子电池，以提供较长的续航时间和较高的能量密度；电控系统将采用基于Raspberry4B作为智能仓储机器人的主控制器，如图\ref{fig:OIP}所示，具有强大的处理能力和丰富的接口，能够满足AGV小车在复杂仓储环境中的控制需求。便于后期维护和升级。\par
           \begin{figure}[htbp]
                \centering
                \includegraphics[width=0.64\textwidth]{OIP.png}
                \caption{树莓派4B示意图}
                \label{fig:OIP}
             \end{figure}
         }
\chapter{AGV机械结构设计}
    \section{核心部件结构设计}
         \xfbody{AGV小车的主要核心部件包括底盘，驱动轮，罩壳，顶升机构等。

         }
    \section{总体结构设计}
         \xfbody{

         } 
    \section{结构强度校核和有限元分析}
             \subsection{力学分析选择}
           \xfbody{通过上述的设计与计算，我们选择对AGV浮动支撑梁进行受力分析。AGV本体重量为 200 kg，额定载荷为 25 kg，因此在满载情况下系统总重量为：\par
           
           \begin{equation}
           W = (200 + 25) \cdot 9.81 \approx 2205 \, \text{N}
           \end{equation}\par   
           由于AGV采用两轮支撑结构，整体重量将由两根支撑梁共同承担。假设受力分布均匀，则单根支撑梁所承受的最大力为总重量的一半，即：\par
           \begin{equation}
           F = \frac{2205}{2} \approx 1102.5 \, \text{N}
           \end{equation}\par
          在实际设计中，为了保证安全性与可靠性，通常会考虑一定的安全系数。若取安全系数 1.02 左右，则单根支撑梁的设计受力可近似为：\par
           \begin{equation} 
           F_{\text{max}} \approx 1125 \, \text{N}
           \end{equation}\par
           因此，在后续的结构设计与强度校核过程中，我们将以 1125 N 作为单根支撑梁的最大受力值进行分析。这一数值将作为梁的截面选择、材料强度校核以及疲劳寿命评估的重要依据。\par
           }
           \subsection{支撑材料选择}
               \xfbody{在材料选择方面，我们不仅考虑了结构强度和安全系数，还综合分析了经济成本、零件的加工工艺特性以及实际使用环境的适应性。经过对比与权衡，最终确定选用 Q235A 普通碳素结构钢 作为支撑梁的材料。该钢材在机械制造与工程结构中应用广泛，具有良好的焊接性能、塑性和韧性，且价格相对低廉，能够满足本设计的承载需求与经济性要求。\par
               其主要性能参数如下表所示：
               \begin{table}[htbp]
                    \wuhao
                    \centering
                    \caption{Q235A号钢的主要性能参数}
                    \begin{tabular}{cc}
                         \toprule \\[-0.5em]
                         性能指标 & 数值范围/说明 \\[0.5em]
                         \midrule \\[-0.5em]
                         屈服强度 & $\geq 235 \, \text{MPa}$ \\[0.5em]
                         抗拉强度 & $370 \sim 500 \, \text{MPa}$ \\[0.5em]
                         延伸率 & $\geq 26\%$ \\[0.5em]
                         弹性模量 & $\approx 2.06 \times 10^{5} \, \text{MPa}$ \\[0.5em]
                         密度 & $7.85 \, \text{g/cm}^3$ \\[0.5em]
                         焊接性能 & 良好 \\[0.5em]
                         冲击韧性 & 常温下具备良好韧性 \\[0.5em]
                         \bottomrule
                    \end{tabular}
               \end{table}
               }
               \subsection{有限元分析}
                   \xfbody{选好材料以后，我们对支撑梁进行有限元分析，以验证其在最大受力情况下的强度和安全性。\par
                   首选选择固定位点如图\ref{fig:fangzhen1}所示\par
                \begin{figure}[htbp]
                \centering
                          \includegraphics[width=0.64\textwidth]{fangzhen1.png}
                \caption{固定支撑梁的有限元分析示意图}
                \label{fig:fangzhen1}
             \end{figure}
             固定好后，我们在轮系固定孔处施加1125N的力F。如图\ref{fig:fangzhen2}所示\par
                \begin{figure}[htbp]
                \centering
                          \includegraphics[width=0.64\textwidth]{fangzhen2.png}
                \caption{施加受力的有限元分析示意图}
                \label{fig:fangzhen2}
             \end{figure}
             在进行有限元分析时，首先需要对结构模型进行网格划分。网格划分的目的是将连续的结构离散化为有限数量的单元，以便通过数值方法求解各节点的应力、应变与位移分布。每个单元之间通过节点相互连接，从而共同反映整体结构的力学特性。如图\ref{fig:fangzhen3}所示，支撑梁被划分为多个有限单元，每个单元的尺寸和形状根据结构的复杂程度和受力情况进行优化，以确保分析结果的准确性和计算效率。\par
                \begin{figure}[htbp]
                \centering
                          \includegraphics[width=0.64\textwidth]{fangzhen3.png}
                \caption{支撑梁的网格划分示意图}
                \label{fig:fangzhen3}
             \end{figure}
             在完成网格划分之后，我们对模型进行了算例运行与受力仿真模拟。通过有限元分析软件对支撑梁在实际载荷工况下进行受力计算，可以得到结构在不同位置的应力分布情况。仿真过程中，载荷条件按照 AGV 的自重与额定载荷进行叠加，并施加在支撑梁的受力点上，同时考虑约束条件与边界条件的合理设置，以保证计算结果的准确性。如图\ref{fig:fangzhen4}所示，仿真结果显示支撑梁在最大受力情况下的应力分布情况，验证了设计的合理性与安全性。\par
               \begin{figure}[htbp]
                \centering
                          \includegraphics[width=0.64\textwidth]{fangzhen4.png}
                \caption{支撑梁受力仿真结果示意图}
                \label{fig:fangzhen4}
             \end{figure}
             在模拟结果中，支撑梁的最大应力值出现在受力集中区域。通过对比材料的屈服强度，可以判断结构是否满足安全要求。本次仿真结果显示，支撑梁的最大屈服力为：\par
             \[\sigma_{\mathrm{max}} = 4.891 \times 10^{6} \, \mathrm{N/m^{2}}\]
            远小于 Q235A 钢的屈服强度 235 MPa，说明支撑梁在设计载荷下具有足够的强度和安全裕度，能够满足AGV小车在仓储环境中的运行需求。\par
            在完成受力分析与仿真计算后，我们进一步对结构的位移变形情况进行了模拟评估。通过有限元分析结果可以看出，在最大载荷工况下，支撑梁的整体变形量极小。具体而言，模型的最大位移变形为：\par
            \[\delta_{\mathrm{max}} = 1.538 \times 10^{-3} \, \text{mm}\]
            这一数值表明结构在受力作用下的刚度表现良好，变形量远低于一般工程设计中允许的位移限值。换句话说，支撑梁在承受设计载荷时几乎没有明显的形变，能够保持整体结构的稳定性与精度要求。\par
            从工程角度来看，位移变形的大小直接影响设备的运行可靠性与使用寿命。若变形过大，可能导致零部件之间的干涉或疲劳损伤；而本次仿真结果显示的最大位移仅为 0.001538 mm，说明设计在刚度与强度方面均满足要求，具有较高的安全裕度和使用可靠性。\par
            在受力仿真分析中，我们不仅关注应力与位移的分布情况，还需要对结构的应变进行评估。如图\ref{fig:fangzhen5}所示。\par
               \begin{figure}[htbp]
                \centering
                          \includegraphics[width=0.64\textwidth]{fangzhen5.png}
                \caption{支撑梁应变分布示意图}
                \label{fig:fangzhen5}
             \end{figure}
            应变是衡量材料在外力作用下产生形变的重要参数，它直接反映了材料的弹性响应能力。通过有限元计算结果可以得到，在最大载荷工况下，结构的最大应变为：\par
            \[\varepsilon_{\mathrm{max}} = 2.268 \times 10^{-5}\]
            这一数值表明，支撑梁在受力作用下的变形处于极小范围，远低于材料的屈服应变。换句话说，结构在该工况下仍处于弹性阶段，没有发生塑性变形，说明设计的刚度与强度均满足要求。\par
            从工程角度来看，最大应变的大小不仅影响结构的安全性，还决定了其长期服役性能。较小的应变意味着结构在反复载荷作用下不易产生疲劳损伤，能够保持稳定的工作状态。因此，本次仿真结果进一步验证了所选材料与结构设计的合理性。\par
            在有限元仿真过程中，除了应力与位移的计算外，应变分析同样是评价结构性能的重要指标。如图\ref{fig:fangzhen6}所示。\par
               \begin{figure}[htbp]
                \centering
                          \includegraphics[width=0.64\textwidth]{fangzhen6.png}
                \caption{支撑梁最大应变分布示意图}
                \label{fig:fangzhen6}
             \end{figure}
            应变反映了材料在外力作用下的相对变形程度，是判断结构是否处于弹性阶段的关键参数。通过本次算例运行结果可以得到，在最大载荷工况下，结构的最大应变为：\par
               \[\varepsilon_{\mathrm{max}} = 2.268 \times 10^{-5}\]
           这一数值远低于材料的屈服应变，说明结构在受力作用下仍处于弹性范围，没有发生塑性变形。换句话说，支撑梁在该工况下的变形完全可逆，卸载后能够恢复原状，保证了整体结构的稳定性与安全性。\par
           从工程角度来看，较小的应变不仅意味着结构刚度良好，还表明其在长期服役过程中不易产生疲劳损伤，能够维持较高的可靠性与使用寿命。因此，该仿真结果进一步验证了所选材料与设计方案的合理性。\par
            综上所述,通过对AGV小车支撑梁的有限元分析，我们验证了设计在最大载荷工况下的强度、刚度和安全性。仿真结果显示，支撑梁的最大应力远低于材料的屈服强度，最大位移变形极小，最大应变也处于弹性范围内，说明设计方案具有足够的安全裕度和可靠性，能够满足AGV小车在仓储环境中的运行需求。同时，这些分析结果为后续的结构优化与改进提供了重要的参考依据。\par


            
         }
\chapter{AGV电气系统设计}
    \section{电气系统总体架构} 
         \xfbody{

         }
    \section{电机，传感器等部件选型与设计}
         \xfbody{

         }
    \section{控制电路设计}
         \xfbody{

         }
\chapter{AGV运动学建模与仿真}
    \section{差速驱动底盘运动学方程推导}
         \xfbody{

         }
    \section{MATLAB/Simulink仿真模型搭建}
         \xfbody{

         }
    \section{典型运动轨迹仿真与分析}
         \xfbody{

         }
\chapter{AGV运动控制算法分析}
    \section{PID控制算原理}
         \xfbody{

         }
    \section{双轮差速PID控制器设计}
         \xfbody{

         }
    \section{基本运动控制实现}
         \xfbody{

         }
     \section{控制效果仿真验证}
             \xfbody{
    
             }
\chapter{结论}
    \section{全文工作总结}
         \xfbody{

         }
    \section{不足和改进方向}
         \xfbody{

         }
    \section{应用前景展望}
         \xfbody{

         }
